\scanpage{109}
\mspage{-10-}

$G_{x0}$, \uline{the intersection} $G'\cap T'$ \uline{finite, and} $\operatorname{rank}G'+\operatorname{rank}T'=r$. \uline{Further} $G'$ \uline{and} $T'$ \uline{are uniquely determined as the commutator subgroup and the identity component of the centre of} $G_{x0}$. (b) \uline{The unipotent elements of} $G_x$ \uline{are all in} $G'$.

Part (b) follows from (a) because $G_{x0}$ contains the unipotent elements of $G_x$ by [8, p. 6-15, Cor. 2]. For the proof of (a) we may imbed $x$ in a maximal torus $T$ and use the above notation. If $y$ in $G_x$ is written $y=u\sigma_wb$ as in 2.4 then the uniqueness in 2.4 implies that $u,\sigma_w$ and $b$ are in $G_x$. By 2.1 and 2.2 we get:

2.8 $G_x$ \uline{is generated by} $T$, \uline{those} $X_\alpha$ \uline{for which} $\alpha(x)=1$, \uline{and those} $\sigma_w$ \uline{for which} $wx=x$. Then $G_{x0}$ is generated by $T$ and the $X_\alpha$ alone because the group so generated is connected and of finite index in $G_x$ (see [8, p. 3-01, Th. 1]). Let $G'$ be the group generated by the $X_\alpha$ alone, and let $T'$ be the identity component of the intersection of the kernels of the roots $\alpha$ such that $\alpha(x)=1$. Then $G'$ is semisimple by [8, p. 17-02, Th. 1], and the other assertions of (a) are soon verified.

2.9 Corollary. \uline{In} 2.7 \uline{every maximal torus containing} $x$ \uline{also contains} $T'$.

For in the above proof $T$ was chosen as an arbitrary torus containing $x$.

2.10 \uline{Remark}. That $G_x$ in 2.7 need not be connected, even if $x$ is regular, is shown by the example: $G=\mathrm{PSL}(2)$, $x=\operatorname{diag}(i,-i)$, $i^2=-1$.

\margin{108}
